\documentclass{article}
\usepackage{amsmath}
\usepackage{amssymb}
\usepackage{array}

\title{Trust Model Formulations}
\author{}
\date{}

\begin{document}

\maketitle

\section{Handle Edge Cases}

\subsection{Zero or Negative Velocity}
\begin{itemize}
    \item \textbf{Scenario}: Leader or reference velocity $v_{\text{ref}} \leq 0$.
    \item \textbf{Formulation}:  
    \begin{itemize}
        \item If $v_{\text{ref}} \leq 0$ and leader acceleration $a_{\text{leader}} < 0$ (braking), but follower acceleration $a_y \geq 0$ (not braking):
        \[
        v_{\text{score\_final}} = 0
        \]
        \item Otherwise:
        \[
        v_{\text{score\_final}} = \max\left(1 - \left| \frac{a_y}{a_{\text{max}}} \right|, 0\right)
        \]
        where $a_{\text{max}}$ is the maximum plausible acceleration (e.g., 10 m/s²).
    \end{itemize}
\end{itemize}

\subsection{Very Small or Negative Distances}
\begin{itemize}
    \item \textbf{Scenario}: Measured distance $d_{\text{measured}} \leq 0$ or $d_{\text{measured}} < d_{\text{critical}}$.
    \item \textbf{Formulation}:
    \begin{itemize}
        \item If $d_{\text{measured}} \leq 0$:
        \[
        d_{\text{score}} = 0
        \]
        \item If $0 < d_{\text{measured}} < d_{\text{critical}}$:
        \[
        d_{\text{score}} = \max\left(1 - \frac{d_{\text{critical}} - d_{\text{measured}}}{d_{\text{critical}}}, 0\right)
        \]
        \item Else:
        \[
        d_{\text{score}} = \text{normal\_d\_score}(d_{\text{measured}})
        \]
        where $\text{normal\_d\_score}$ is the standard distance scoring function.
    \end{itemize}
\end{itemize}

\subsection{Extreme Accelerations}
\begin{itemize}
    \item \textbf{Scenario}: Reported acceleration $|a_y| > a_{\text{max}}$.
    \item \textbf{Formulation}:
    \[
    a_{\text{score}} = 
    \begin{cases} 
    0 & \text{if } |a_y| > a_{\text{max}} \\
    \text{normal\_a\_score}(a_y) & \text{otherwise}
    \end{cases}
    \]
    where $\text{normal\_a\_score}$ is the standard acceleration scoring function.
\end{itemize}

\subsection{Beacon Loss}
\begin{itemize}
    \item \textbf{Scenario}: No beacons received for time $t$ since the last beacon.
    \item \textbf{Formulation}:
    \[
    \text{beacon\_score} = e^{-\lambda t}
    \]
    where $\lambda$ is a decay rate (e.g., 0.1). If $t > t_{\text{max}}$ (e.g., 5 seconds):
    \[
    \text{beacon\_score} = 0
    \]
\end{itemize}

\section{Safety-First Design}

\subsection{Asymmetric Trust Adjustment}
\begin{itemize}
    \item \textbf{Concept}: Decrease trust quickly when it drops, increase it slowly when it rises.
    \item \textbf{Formulation}:  
    Let $\text{trust}_{\text{ext}}$ be the current trust, $\text{trust}_{\text{new}}$ the new trust sample, $\alpha_{\text{decrease}} = 0.5$, and $\alpha_{\text{increase}} = 0.1$:
    \[
    \text{trust}_{\text{ext}} = 
    \begin{cases}
    (1 - \alpha_{\text{decrease}}) \cdot \text{trust}_{\text{ext}} + \alpha_{\text{decrease}} \cdot \text{trust}_{\text{new}} & \text{if } \text{trust}_{\text{new}} < \text{trust}_{\text{ext}} \\
    (1 - \alpha_{\text{increase}}) \cdot \text{trust}_{\text{ext}} + \alpha_{\text{increase}} \cdot \text{trust}_{\text{new}} & \text{otherwise}
    \end{cases}
    \]
\end{itemize}

\subsection{Critical Thresholds for Emergency Actions}
\begin{itemize}
    \item \textbf{Concept}: Trigger safety actions if trust falls too low.
    \item \textbf{Formulation}:  
    Let $\text{final\_score}$ be the overall trust score and $\tau_{\text{critical}} = 0.2$:
    \[
    \text{emergency\_action} = 
    \begin{cases} 
    1 & \text{if } \text{final\_score} < \tau_{\text{critical}} \\
    0 & \text{otherwise}
    \end{cases}
    \]
\end{itemize}

\subsection{Conservative Trust Combination}
\begin{itemize}
    \item \textbf{Concept}: Combine local and global trust conservatively.
    \item \textbf{Formulation}:  
    Let $\text{trust}_{\text{sample}}$ be the sample trust, $\gamma_{\text{cross}}$ the cross-verified trust, and $\gamma_{\text{local}}$ the local trust:
    \begin{itemize}
        \item Using minimum:
        \[
        \text{trust}_{\text{ext}} = \min(\text{trust}_{\text{sample}}, \gamma_{\text{cross}} \cdot \gamma_{\text{local}})
        \]
        \item Using product:
        \[
        \text{trust}_{\text{ext}} = \text{trust}_{\text{sample}} \cdot (\gamma_{\text{cross}} \cdot \gamma_{\text{local}})
        \]
    \end{itemize}
\end{itemize}

\section{Monitor Sudden Changes}

\subsection{Basic Change Detection for $\gamma_{\text{cross}}$}
\begin{itemize}
    \item \textbf{Approach}: Detect anomalies using a sliding window of size $w$ (e.g., 10).
    \item \textbf{Formulation}:  
    Let $\gamma_{\text{cross}}^{(t)}$ be the value at time $t$, and $\Gamma = [\gamma_{\text{cross}}^{(t-w+1)}, \dots, \gamma_{\text{cross}}^{(t)}]$:
    \[
    \mu_{\gamma} = \frac{1}{w} \sum_{i=1}^{w} \Gamma_i, \quad \sigma_{\gamma} = \sqrt{\frac{1}{w} \sum_{i=1}^{w} (\Gamma_i - \mu_{\gamma})^2}
    \]
    If $|\gamma_{\text{cross}}^{(t+1)} - \mu_{\gamma}| > 2 \sigma_{\gamma}$:
    \[
    \text{anomaly}_{\gamma}^{(t+1)} = 1
    \]
    Else:
    \[
    \text{anomaly}_{\gamma}^{(t+1)} = 0
    \]
\end{itemize}

\subsection{Monitoring State Components (e.g., Position, Velocity)}
\begin{itemize}
    \item \textbf{Concept}: Check for sudden changes in differences between target and host states.
    \item \textbf{Formulation}:  
    For a state component $s$ (e.g., position $x$, velocity $v$):
    \[
    s_{\text{diff}}^{(t)} = s_{\text{target}}^{(t)} - s_{\text{host}}^{(t)}
    \]
    Using a sliding window $S = [s_{\text{diff}}^{(t-w+1)}, \dots, s_{\text{diff}}^{(t)}]$:
    \[
    \mu_s = \frac{1}{w} \sum_{i=1}^{w} S_i, \quad \sigma_s = \sqrt{\frac{1}{w} \sum_{i=1}^{w} (S_i - \mu_s)^2}
    \]
    If $|s_{\text{diff}}^{(t+1)} - \mu_s| > 2 \sigma_s$:
    \[
    \text{anomaly}_s^{(t+1)} = 1
    \]
    Else:
    \[
    \text{anomaly}_s^{(t+1)} = 0
    \]
\end{itemize}

\section{Log Flags Separately}

\begin{itemize}
    \item \textbf{Concept}: Record anomalies without immediately altering trust.
    \item \textbf{Formulation}:  
    At each time step $t$:
    \begin{itemize}
        \item Store trust:
        \[
        \text{trust\_sample\_log}[t] = \text{trust}_{\text{ext}}^{(t)}
        \]
        \item Store anomaly flags:
        \[
        \text{anomaly\_flags}[t] = 
        \begin{bmatrix}
        \text{anomaly}_{\gamma}^{(t)} \\
        \text{anomaly}_{\text{position}}^{(t)} \\
        \text{anomaly}_{\text{velocity}}^{(t)} \\
        \vdots
        \end{bmatrix}
        \]
        \item Optional cumulative check:
        \[
        \text{anomaly\_count} = \sum_{i=t-w+1}^{t} \text{anomaly}_{\gamma}^{(i)}
        \]
        If $\text{anomaly\_count} > k$ (e.g., 3):
        \[
        \text{trust}_{\text{ext}}^{(t)} = \text{trust}_{\text{ext}}^{(t)} \cdot \beta
        \]
        where $\beta = 0.5$ reduces trust.
    \end{itemize}
\end{itemize}



\section{Combining Local and Global Trust into a Single Trust Score}
You want to combine local trust (based on direct observations like sensor data) and global trust (based on platoon-wide consistency) into one trust score for each vehicle. This makes sense because a single score simplifies the design of the observer gains and communication across the platoon. A straightforward way to do this is to multiply the two trust values:

\[
\text{trust}_{i \to j} = \text{trust}_{\text{local}, i \to j} \cdot \text{trust}_{\text{global}, i \to j}
\]

\begin{itemize}
    \item \textbf{Why this works}: Multiplying ensures that if either local or global trust is low, the combined score reflects that distrust. For example, if a nearby vehicle has faulty sensors (low local trust) or doesn't align with the platoon's state (low global trust), its overall trust score will be low. This conservative approach prioritizes safety.
    
    \item \textbf{How it's used}: Each vehicle \( i \) calculates \( \text{trust}_{i \to j} \) for every vehicle \( j \) it observes and sends this to the host vehicle for aggregation.
\end{itemize}

\section{Logging Flags Separately Without Altering Trust Immediately}
You also want to log flags for nearby vehicles separately, without letting them instantly change the trust score. This is a smart move because it gives you flexibility: you can monitor issues over time without overreacting to temporary glitches.

\begin{itemize}
    \item \textbf{What are flags?}: These are indicators of specific anomalies, like:
    \begin{itemize}
        \item \( \text{flag}_{\text{position}, j} \): Unusual position data.
        \item \( \text{flag}_{\text{velocity}, j} \): Unexpected velocity changes.
        \item \( \text{flag}_{\text{global}, j} \): Inconsistencies with platoon-wide trust.
    \end{itemize}
    
    \item \textbf{How to log them}: In the host vehicle, maintain a log (e.g., a table or list) for nearby vehicles. For each time step, record any flags triggered by a vehicle \( j \) without modifying \( \text{trust}_{i \to j} \) right away.
    
    \item \textbf{Why this helps}: Logging flags separately lets you:
    \begin{itemize}
        \item Track patterns (e.g., a vehicle consistently showing position errors).
        \item Avoid knee-jerk reactions to one-off issues (e.g., a single noisy measurement).
    \end{itemize}
\end{itemize}

For example, your log might look like this over a few time steps for a nearby vehicle \( j \):

\begin{center}
    \begin{tabular}{cccc}
        \toprule
        Time Step & Position Flag & Velocity Flag & Global Flag \\
        \midrule
        1         & 0             & 0             & 0           \\
        2         & 1             & 0             & 0           \\
        3         & 1             & 1             & 0           \\
        \bottomrule
    \end{tabular}
\end{center}

Here, 1 means an anomaly was detected, and 0 means no issue.

\section{Using Flags for Nearby Vehicles When Designing Observer Gains}
When designing the observer gains (e.g., weights that determine how much you trust each vehicle's data), you want to check if the target vehicle \( j \) being estimated is nearby. If it is, you can use the logged flags to adjust or prioritize the trust score before aggregating it with trust from other vehicles.

\subsection{General Observer Gain Design}
First, let's define the basic weights without flags:
\begin{itemize}
    \item \textbf{Consensus weight \( w_{il}^{(j)} \)}: How much vehicle \( i \) trusts vehicle \( l \)'s estimate of vehicle \( j \)'s state:
    \[
    w_{il}^{(j)} = \frac{\text{trust}_{i \to l}}{\sum_{m \in N_i} \text{trust}_{i \to m}}
    \]
    This normalizes the trust across all neighbors \( N_i \).
    
    \item \textbf{Tracking weight \( w_{i0}^{(j)} \)}: How much vehicle \( i \) trusts its own measurement of \( j \):
    \[
    w_{i0}^{(j)} = \text{trust}_{i \to j}
    \]
\end{itemize}

These weights use the combined trust score directly and are fine for distant vehicles.

\subsection{Adjusting for Nearby Vehicles Using Flags}
For nearby vehicles, you can add an extra layer of caution by using the flags. The idea is to penalize the trust score of a nearby vehicle if its log shows frequent anomalies, prioritizing safety since these vehicles are more critical (e.g., for collision avoidance).

\begin{itemize}
    \item \textbf{Step 1: Calculate an anomaly fraction}:
    \begin{itemize}
        \item For a nearby vehicle \( l \), look at the log over the last \( T \) time steps (e.g., \( T = 10 \)).
        \item Count how often any flag (position, velocity, or global) was triggered.
        \item Compute the fraction \( f_l \):
        \[
        f_l = \frac{\text{Number of time steps with any flag}}{\text{Total time steps } T}
        \]
        For example, if vehicle \( l \) had flags in 3 out of 10 steps, \( f_l = 0.3 \).
    \end{itemize}
    
    \item \textbf{Step 2: Apply a penalty to the trust score}:
    \begin{itemize}
        \item Reduce the trust score based on \( f_l \):
        \[
        \text{trust}_{i \to l}^{\text{adjusted}} = \text{trust}_{i \to l} \cdot (1 - \alpha f_l)
        \]
        where \( \alpha \) is a tuning parameter (e.g., 0.5) that controls how much the flags lower the trust. If \( \alpha = 0.5 \) and \( f_l = 0.3 \):
        \[
        \text{trust}_{i \to l}^{\text{adjusted}} = \text{trust}_{i \to l} \cdot (1 - 0.5 \cdot 0.3) = \text{trust}_{i \to l} \cdot 0.85
        \]
    \end{itemize}
    
    \item \textbf{Step 3: Use the adjusted trust for nearby vehicles}:
    \begin{itemize}
        \item For a nearby vehicle \( l \), replace \( \text{trust}_{i \to l} \) with \( \text{trust}_{i \to l}^{\text{adjusted}} \) in the weight calculation:
        \[
        w_{il}^{(j)} = \frac{\text{trust}_{i \to l}^{\text{adjusted}}}{\sum_{m \in N_i} \text{trust}_{i \to m}^{\text{adjusted}}}
        \]
        \item This reduces the influence of a nearby vehicle with frequent anomalies.
    \end{itemize}
\end{itemize}

\subsection{Why This Prioritizes Nearby Vehicles}
By checking if the target vehicle is nearby and adjusting its trust based on flags, you ensure that:
\begin{itemize}
    \item Vehicles with persistent issues (e.g., faulty sensors) have less influence on your estimate, protecting the host vehicle.
    \item The adjustment only applies to nearby vehicles, where safety is most critical, while distant vehicles use the unadjusted trust score.
\end{itemize}

\section{Aggregating Trust Across the Platoon}
After calculating trust scores (adjusted for nearby vehicles), the host vehicle aggregates them to get a platoon-wide trust for each vehicle \( j \):
\[
\text{trust}_{\text{agg}, j} = \frac{1}{m} \sum_{i=1}^{m} \text{trust}_{i \to j}
\]
where \( m \) is the number of vehicles reporting trust for \( j \). For nearby vehicles, the host uses the adjusted trust scores based on its own logs.

\begin{itemize}
    \item \textbf{Host's special role}: If \( j \) is nearby, the host can further tweak \( \text{trust}_{\text{agg}, j} \) using its own logged flags:
    \[
    \text{trust}_{\text{host}, j} = \text{trust}_{\text{agg}, j} \cdot (1 - \beta f_j^{\text{host}})
    \]
    where \( f_j^{\text{host}} \) is the anomaly fraction from the host's log, and \( \beta \) (e.g., 0.5) is a tuning parameter.
\end{itemize}

This ensures the host prioritizes its own observations for nearby vehicles while still considering the platoon's input.

\section{Putting It All Together}
Here's how it works in the host vehicle:
\begin{enumerate}
    \item \textbf{Compute combined trust}: For every vehicle \( j \), calculate \( \text{trust}_{i \to j} = \text{trust}_{\text{local}, i \to j} \cdot \text{trust}_{\text{global}, i \to j} \).
    \item \textbf{Log flags}: For nearby vehicles, record anomalies in a log without changing the trust score immediately.
    \item \textbf{Design gains}:
    \begin{itemize}
        \item For distant vehicles, use \( \text{trust}_{i \to j} \) directly to compute weights.
        \item For nearby vehicles, check the log, compute \( f_l \), adjust the trust to \( \text{trust}_{i \to l}^{\text{adjusted}} \), and use that for weights.
    \end{itemize}
    \item \textbf{Aggregate trust}: Combine trust scores across the platoon, applying host-specific adjustments for nearby vehicles based on its own flags.
\end{enumerate}

\section{Why This Approach Makes Sense}
\begin{itemize}
    \item \textbf{Simplicity}: A single trust score keeps the observer design and communication straightforward.
    \item \textbf{Safety}: Logging flags and adjusting trust for nearby vehicles ensures you're cautious where it matters most.
    \item \textbf{Flexibility}: Flags don't alter trust instantly, so you can decide how to use them—whether for real-time tweaks or later analysis.
\end{itemize}


\end{document}